\section{Vector Spaces}

\begin{definition}
$n$-dimensional Affine Space (Vector Space) $\R_n$. An ordered $n$-tuple of real numbers, i.e., a system of $(a_1, a_2, ..., a_n)$ of $n$ real numbers , called components, in a definite order, is called a point ($n$ a positive integer). The numbers $a_1,a_2, ...,a_n$ are called the coordinates of the point; in particular $a_1$ is called its first, $a_2$ its second, ..., $a_n$ its $n$-th coordinate. Two points $P=(a_1,a_2,...,a_n)$ and $Q = (b_1,b_2,...,b_n)$ are said to 
be the {\elevenit same} or to coincide if and only if $a_1 = b_1, a_2 = b_2, ..., a_n = b_n$. The totality of all $n$-tuples of real numbers is called $n$-dimensional Affine Space and is denoted by $\R_n$.
\end{definition}

\begin{definition}[Vector Notation]
An element, $\x \in \R_n$ is defined by its $n$ components $(x_1, x_2, \hdots, x_n)$, as follows
\[ \x = (x_1, x_2, \hdots, x_n) \]
\end{definition}

\begin{definition}[Linear Dependence/Independence]
In the abstract vector space $V$ the finite set of vectors
$\x_1, \x_2, \hdots, \x_n$ is said to be linear dependent if scalars $\lambda_1, \lambda_2, \hdots, \lambda_n$ exist, not all zero, such that 
\[ \lambda_1 \x_1 + \lambda_2 \x_2 + \hdots + \lambda_n\x_n = \o.\] If no such scalars exist, i.e., if 
\[ \lambda_1 \x_1 + \lambda_2 \x_2 + \hdots + \lambda_n\x_n = \o \implies \lambda_i  = 0 ~\forall ~i = 1,2, \hdots, n\] then the set is linearly independent.  
\end{definition}

{\bf Conventions}: We adopt the Range and Summation Conventions:
\begin{definition}
{\bf Range Convention.} {\it When a small Latin suffix (superscript or subscript) occurs unrepeated in a term, it is understood to take all the values
$1,2,...,n$, where n is the number of dimensions of the space.}
\end{definition}

\begin{definition}
{\bf Summation Convention.} {\it When a small Latin suffix is repeated in a term, summation with respect to that suffix is understood, the range of summation
being $1,2,...,n$.}
\end{definition}

{\bf Kronecker delta}:
\[ \delta^s_r = \begin{cases} 1 & \mbox{if } s = r \\0 & \mbox{if } r \ne s\end{cases} \]
Also, the another form of the Kronecker delta is $\delta_{rs}$ which satisfies the same cases as above. \\

{\bf Permutation Symbol}: Consider a set $\sigma$ of $\sigma_1, \sigma_2, ..., \sigma_n$ numbers taken from the set $\{1,2,...,n\}$. Let
\[ \epsilon_{\sigma_1\sigma_2...\sigma_n} =
\begin{cases}
+1 & \mbox{if }\sigma \mbox{ is an even permutation of $\{1, 2, ..., n\}$}\\
-1 & \mbox{if }\sigma \mbox{ is an odd permutation of $\{1, 2, ..., n\}$}\\
0 & \mbox{if }\sigma \mbox{ has any duplicates}
\end{cases} \]

\begin{definition}[Span]
The set of vectors $\x_1, \x_2, \hdots, \x_n$ are said to span or generate, the space $V$ if any vector $\x$ of $V$ is expressible as a linear combination of them, i.e. if 
\[ \x = \lambda_1 \x_1 + \lambda_2 \x_2 + \hdots + \lambda_n \x_n \] for some scalars $\lambda_1, \lambda_2, \hdots, \lambda_n$ not necessarily unique. 
\end{definition}

\begin{definition}[Basis]
The set $\e_1, \e_2, \hdots, \e_n$ is a Basis of $V$: 1) if they are linearly independent; and (2) they span $V$. 
\end{definition}

\begin{theorem}
If $\e_1, \e_2, \hdots, \e_n$ span $V$ then they form a basis if and only if $\x$ in $V$ is expressible uniquely as a linear combination of $\e_1, \e_2, \hdots, \e_n$ .
\end{theorem}

\begin{theorem}
If $\x_1, \x_2, \hdots, \x_m$ span $V$, then there is a subset of these which is a basis of $V$. \label{SpanningSubset}
\end{theorem}
\begin{proof}
If $\x_1, \x_2, \hdots, \x_m$ are linearly independent, they form a basis and the theorem is proved. If they are linearly dependent then some $\x_{k+1}$ is a linear
combination of 
\[  \x_1, x_2, \hdots, \x_k  \mbox{ say } \x_{k+1} = \sum_1^k\, \mu_i \x_i ,  \]
Consider the set $\x_1, \x_2, \hdots, \x_k, \x_{k_2} \hdots, \x_m$. Any vector $\x$ can be expressed as $\displaystyle \x = \sum_1^m \, \lambda_i \x_i$, since the $\x_i$'s also span $V$. 
Hence  
\[ \x = \sum_{\substack{i = 1\\ i\ne k+1j}} \, \lambda_i \x_i  +\lambda_{k+1} \sum_1^k\, \mu_i\x_i  \]
and is a linear combination of $\x_1, \x_2, \hdots, \x_k, \x_{k+2}, \hdots, \x_m$ so that these also span $V$.\\

If this set is now linearly independent the theorem is proved as they form a basis. If not repeat the above argument and obtain a spanning set of $m-2$ vectors, which may or may not
be linearly independent. Repeat until after a most $m-1$ sets we do obtain a basis, which is a subset of the original set. This must occur since after $m-1$ steps we are left with a single vector
which spans $V$, and this must be non-zero unless $V$ is the trivial space $\{ \o \}$ and hence is linearly independent, since $\lambda_1 \x_1 = \o$ with $\x_1 \ne \o \implies \lambda_1 = 0$. 
\end{proof}

\begin{corollary}
If $V$ possesses a finite spanning set then $V$ is finite dimensional. 
\end{corollary}

\begin{theorem}
If $\x_1, \x_2, \hdots, \x_n$ are linearly independent and $V$ is finite dimensional then there exists a basis containing
\[ \x_1, \x_2, \hdots, \x_p.\]
\end{theorem}
\begin{proof}
Let $\e_1, \e_2, \hdots, \e_n$ be a basis of $V$. Consider the set $\x_1, \x_2, \hdots, \x_p, \e_1, \hdots, \e_n$. This set certainly spans $V$. Using the argument of Thm. (\ref{SpanningSubset}) on this set, but notice that 
at each stage we obtain one of the $\e_i$'s as a linear combination of the preceding vectors, since no $\x$ can be a linear combination of the preceding $\x_i$'s as the $\x_i$'s form a linearly independent set. \\

Hence we are left with the $\x_i$'s and the process must stop before or exactly when the $\e_i$'s are exhausted. We this obtain a basis which contains $\x_1, \x_2, \hdots, \x_p$ together with some of the $\e_i$'s. 
\end{proof}

\begin{theorem}[Steinitz Exchange Lemma]
Let $U$ and $W$ be finite subsets of a vector space $V$ and let $U$ be a set of linearly independent vectors and let $W$ span $V$. 
If $\{\x_1, \x_2, \hdots, \x_p\} \in U$ is a linearly independent set, and $\{\y_1, \y_2, \hdots, \y_m\} \in W$ spans $V$, then $p\le m$. \label{SubBasis}
\end{theorem}

\begin{proof}
Consider the set $\x_p, \y_1, \y_2, \hdots, \y_m$, Since the $\y_i$'s span $V$, $\x_p$ is a linear combination of them and so the set is linearly dependent. Hence some $\y_{k+1}$ is a linear combination of 
\[ x_p, \y_1, \hdots, \y_k.\] This means that $\x_p, \y_1, \hdots, \y_k, \y_{k+2}, \hdots, \y_m$ span $V$. \\

Continuing similarly, adding $\x_{p-1}$ to this set. We obtain a linearly dependent set and some $\y_{l_1}$ is expressible ad a linear combination of the others, so that we have a spanning set comprising
$\x_{p-1}, \x_p,$ and all the $\y_i$'s except $\y_{k+1}$ and $\y_{l + 1}$. By continuing this process we obtain a spanning set consisting of one more $\x$ and one less $\y$. We always obtain a $\y_i$ since the $\x_i$'s are 
linearly independent and so no $\x_{i+1}$ is a linear combination of the preceding $\x_i$'s that have been introduced. \\

The process must be possible until all the $\x_i$'s have been exhausted and we must at each stage have some $\y_i$'s left, otherwise the set consisting of the $\x_i$'s along could not be linearly independent. 
Hence we finally obtain a spanning set $\x_1, x_2, \hdots, \x_p$ together with $(m-p)\y_i$'s and thus $m-p \ge 0$. Therefore $p\le m$. 
\end{proof}

\begin{theorem}
If $\e_1, \e_2, \hdots, \e_n$ and $\f_1, \f_2, \hdots, f_m$ are both basis for $V$, then $m = n$. 
\end{theorem}

\begin{proof}
Both sets are linearly independent and both span $V$. Using Thm. (\ref{SubBasis}), since the $\e_i$'s are linearly independent and the $\f_i$'s span $V$,we have $n\le m$. Since the $\e_i$'s span $V$ and the $\f_i$'s are linearly 
independent, $m\le n$. Therefore $m = n$. 
\end{proof}

If $V$ possesses a basis, then all the same number of elements. We say that $V$ is {\twelveit finite dimensional} and the number of elements in the basis is called the dimension of $V$ and is denoted by $\dim V$. If 
$V$ possesses no finite basis, then we say that $V$ is {\twelveit infinite dimensional}. 

\begin{theorem}
If $V$ has dimension $n$, there cannot exist a set of more than $n$ linearly independent vectors in $V$. 
\end{theorem}

\begin{theorem}
If $V$ has dimension $n$, there cannot exist a set of less than $n$ vectors which span $V$. 
\end{theorem}
